% Notes of the Montreal seminar, 1970 (cote 129, typed part): archive cote 129, site page 9 (PDF page 10).
% typescript. Draft compared: none.
% AI-transcribed by Claude Opus 5.5 (claude-opus-5-5, adaptive thinking, effort high), Message Batches API, batch mtl129_full (msgbatch_01AGnDKHwy6ATCM7ehUa6263), 2026-11-13; not yet checked by eye.
\scanpage{129/9}
en faisant opérer $\lambda \in W$ sur $(\mathbb{W}_n)_k$ par multiplication par $F^{1-n}(\lambda)$.

Remarquons que l'élément $F^{1-n}(\lambda) \in W$ est bien défini car,
$k$ étant parfait, $F\colon W \to W$ est un isomorphisme. Cette multiplication
a un sens car, pour toute $k$-algèbre $A$, $\mathbb{W}_n(A)$ est une $W$-algèbre via le
morphisme
\[ W = \mathbb{W}(k) \xrightarrow{\ \pi_n\ } \mathbb{W}_n(k) \to \mathbb{W}_n(A) . \]

Pour que cette action se transporte à $\underrightarrow{\mathbb{W}}_k$, il faut montrer que
l'action de $\lambda \in W$ commute aux morphismes $T$. On vérifie immédiatement
qu'il suffit de montrer que pour toute $k$-algèbre $A$, on a:
\[ \forall\, x \in \mathbb{W}_n(A) ,\ \forall\, \lambda' \in \mathbb{W}_{n+1}(A) \quad \lambda'.T(x) = T(FR(\lambda').x) \]
ou, ce qui revient au même puisque $R$ est une surjection:
\[ \forall\, \lambda' ,\ y \in \mathbb{W}_{n+1}(A) ,\ \ \lambda'.V(y) = V(F(\lambda'),y) . \]
Or si $A$ est un anneau parfait, $F$ est un isomorphisme, et cette relation
est vérifiée car
\[ F(\lambda'.V(y)) = F(\lambda')\ FV(y) = F(\lambda').py = p[F(\lambda').y] = FV[F(\lambda').y]. \]
Dans le cas général, il suffit de remarquer que $\lambda'.V(y)$ et $V(F(\lambda').y)$
s'expriment à l'aide de polynômes à coefficients dans $\mathbb{Z}/p\mathbb{Z}$. D'après
ce qui précède, ces polynômes prennent la même valeur si les variables
sont prises dans un anneau parfait, par exemple la clôture algébrique
de $k$. On en déduit qu'ils sont égaux.

\mspage{\hand{9}}
